\scanpage{41}
On voudrait trouver des foncteurs (adjoints ?) en sens inverse
\diagram{\begin{tikzcd}[ampersand replacement=\&,row sep=small]\unsure{f_n^\circ|\mathfrak f_n^\circ}\&G_n\text{-ens}\arrow[l]\arrow[dd,dashed]\\\rotatebox{90}{$\supset$}\&{}\\f_n\&(\mathrm{Casc})_n\arrow[l,"\psi_n"]\end{tikzcd}}
On définira $\psi_n$ par récurrence sur $n$.

\textcircled{1} $n=0$ \qquad $\psi_0(E_0)=E_0$ avec top. discr.
\textcircled{2} Cas général.
\[ \overbrace{E_n\to\underbrace{\overbrace{E_{n-1}\to\ldots\to E_1\to E_0}^{C_{n-1}}}_{\substack{\text{définit }X_{n-1}\supset X_{n-2}\supset\ldots\\\text{par hyp. de réc.}\\X_{n-1}=\unsure{\varphi_{n-1}|\psi_{n-1}}(C_{n-1})}}}^{C_n} \]
On aura \quad $\pi_0(\widetilde X_n)\simeq E_n/G_{n-1}$
$\partial\widetilde X_n$ est défini par \deleted{\illegible{1}} ($E_n$ avec action de $G_{n-1}$) par l'hyp. de récurrence
\[ \begin{aligned} \deleted{\unsure{\partial}}X_{n-1}&=\psi_{n-1}(C_{n-1})\\ \partial\widetilde X_n&=\psi_{n-1}(\rho_{n-1}(E_n)) \end{aligned} \]
$\partial\widetilde X_n\to X_{n-1}$ déduit en appliquant $\psi_{n-1}$ à
\[ \rho_{n-1}(E_n)\longrightarrow C_{n-1} \]
\margin{préciser ce point}

Donc on définira
\[ \psi_n(C_n)=\text{\deleted{\unsure{Cône}} \inserted{Multicône} de }\psi_{n-1}\bigl(\rho_{n-1}(E_n)\to C_{n-1}\bigr) \]
\textemdash\,\textemdash\,\textemdash\quad Il faudrait donc expliciter en un \uline{Th 4} Description de $\psi_n\colon(\mathrm{Casc})_n\to\mathcal F_n$, et le cas échéant propriété d'adjonction.
% archive sheet number: 40
